% --- Archon physics formalization source begin ---
% archon:physics
% archon:covers PhyXMiniProblems/problem_phyx_mini_0610.lean
% archon:source-report reports/phyx_mini/problem_phyx_mini_0610.source.json

\chapter{Physics problem phyx\_mini\_0610}
\label{ch:PhyXMiniProblems_problem_phyx_mini_0610}

\paragraph{Problem source.}
\#\# Physical scenario

A fluorescent tube with length \$L\$ is fixed horizontally above a ticket window in a train station. At time \$t = 0\$ the tube lights up, changing color from white to bright yellow. After a duration \$T\$ the tube turns off, reverting to its white color. The tube lies along the \$x\$-axis of frame \$S\$ fixed with respect to the station, with its right end at the origin. We can represent the history of the tube using a “spacetime” diagram, as shown in figure. Events 1 and 2 correspond, respectively, to the right and left ends of the tube at time \$t = 0\$, while events 3 and 4 correspond to the right and left ends of the tube at time \$t = T\$. The totality of events corresponding to positions on the tube when it is lit correspond to the yellow region on the diagram. The tube is observed from the window of a rocket train passing the station from left to right at speed \$v\$. The frame \$S'\$ is fixed with respect to the train such that the \$x'\$-axis coincides with the \$x\$-axis of frame \$S'\$, and such that event 1 occurs at the origin \$(x', t') = (0, 0)\$.We use the product \$ct\$ rather than \$t\$ on the vertical axis of our figure so that both axes have the dimension of distance. Thus the yellow region has a well-defined area.

\#\# Question

What is the area of the yellow region in the \$S\$ frame plot?

\#\# Answer choices

- A: A: cLT

- B: B: c/LT

- C: C: 2cLT

- D: D: cL/T

\#\# Figure caption (auxiliary; use the image as primary evidence)

The image is a diagram featuring a square with four vertices labeled 1, 2, 3, and 4. The square is shaded in yellow and lies in a coordinate system with an x-axis horizontally and a vertical axis labeled "ct." 

- Vertex 1 is at the origin.
- Vertex 2 is on the x-axis labeled as \textbackslash{}(-L\textbackslash{}).
- Vertex 3 is on the vertical axis labeled as \textbackslash{}(cT\textbackslash{}).
- Vertex 4 is positioned at the intersection of vertical line \textbackslash{}(x = -L\textbackslash{}) and horizontal line \textbackslash{}(ct = cT\textbackslash{}).

The square's edges are drawn with solid lines for the bottom and right sides and dashed lines for the top and left sides.

\#\# Dataset metadata

Relativity; Physical Model Grounding Reasoning, Spatial Relation Reasoning

\paragraph{Recorded answer/context.}
A: A: cLT

\paragraph{Figure/image path.}
/root/proposal\_for\_physic/science-mango/phyx\_mini\_run/phyx\_data/test\_image/610.png

\paragraph{Formalization target.}
create a compiling Lean file with sorry bodies at `PhyXMiniProblems/problem\_phyx\_mini\_0610.lean`. The Lean declarations must preserve the physical quantities, dimensions or dimensional roles, figure labels, governing-law hypotheses, and final relation expressed by this problem.
Use Mathlib/Physlib names found through LeanExplore where available. If a physics API is missing, introduce faithful local abstractions rather than scalar placeholder aliases.

\begin{theorem}[Physics formalization target]
\label{thm:physics:phyx_mini_0610:target}
\lean{PhyXMiniProblems.ProblemPhyXMini0610.problem_phyx_mini_0610}
\uses{def:physics:phyx-mini-0610:phyxminiproblems-problemphyxmini0610-lengthreadout, def:physics:phyx-mini-0610:phyxminiproblems-problemphyxmini0610-timereadout, def:physics:phyx-mini-0610:phyxminiproblems-problemphyxmini0610-speedreadout, def:physics:phyx-mini-0610:phyxminiproblems-problemphyxmini0610-areareadout, def:physics:phyx-mini-0610:phyxminiproblems-problemphyxmini0610-fluorescenttubespacetimesetup, def:physics:phyx-mini-0610:phyxminiproblems-problemphyxmini0610-matchesfluorescenttubescenario, def:physics:phyx-mini-0610:phyxminiproblems-problemphyxmini0610-matchesprimaryspacetimefigure, def:physics:phyx-mini-0610:phyxminiproblems-problemphyxmini0610-matcheseventcoordinatedata, def:physics:phyx-mini-0610:phyxminiproblems-problemphyxmini0610-hasphysicalparameters, def:physics:phyx-mini-0610:phyxminiproblems-problemphyxmini0610-satisfieslightscaledtimecoordinate, def:physics:phyx-mini-0610:phyxminiproblems-problemphyxmini0610-satisfiesilluminatedtubeworldsheet, def:physics:phyx-mini-0610:phyxminiproblems-problemphyxmini0610-satisfiesaxisalignedrectanglearealaw}
The assigned autoformalize agent should translate this physics problem into Lean declarations in the covered file, with theorem and lemma proof bodies written as `by sorry`.
\end{theorem}
\begin{proof}
This is an autoformalization task, not a proof task. Produce statements that can later be proved by the physics prover without weakening the physical model.
\end{proof}

% --- Archon declaration topology begin ---
\section{Declaration topology}
\begin{definition}[Length Quantity]
  \label{def:physics:phyx-mini-0610:phyxminiproblems-problemphyxmini0610-lengthquantity}
  \lean{PhyXMiniProblems.ProblemPhyXMini0610.LengthQuantity}
  A nonnegative, unit-independent physical tube length
\end{definition}
\begin{proof}
  This is the declared model definition of Length Quantity.
\end{proof}
\begin{definition}[Signed Length Quantity]
  \label{def:physics:phyx-mini-0610:phyxminiproblems-problemphyxmini0610-signedlengthquantity}
  \lean{PhyXMiniProblems.ProblemPhyXMini0610.SignedLengthQuantity}
  A signed, unit-independent spatial or ct coordinate
\end{definition}
\begin{proof}
  This is the declared model definition of Signed Length Quantity.
\end{proof}
\begin{definition}[Time Quantity]
  \label{def:physics:phyx-mini-0610:phyxminiproblems-problemphyxmini0610-timequantity}
  \lean{PhyXMiniProblems.ProblemPhyXMini0610.TimeQuantity}
  A nonnegative, unit-independent physical duration
\end{definition}
\begin{proof}
  This is the declared model definition of Time Quantity.
\end{proof}
\begin{definition}[Signed Speed Quantity]
  \label{def:physics:phyx-mini-0610:phyxminiproblems-problemphyxmini0610-signedspeedquantity}
  \lean{PhyXMiniProblems.ProblemPhyXMini0610.SignedSpeedQuantity}
  A signed one-dimensional physical velocity, including v and c
\end{definition}
\begin{proof}
  This is the declared model definition of Signed Speed Quantity.
\end{proof}
\begin{definition}[length Readout]
  \label{def:physics:phyx-mini-0610:phyxminiproblems-problemphyxmini0610-lengthreadout}
  \lean{PhyXMiniProblems.ProblemPhyXMini0610.lengthReadout}
  \uses{def:physics:phyx-mini-0610:phyxminiproblems-problemphyxmini0610-lengthquantity}
  Read a nonnegative physical length in a selected length unit
\end{definition}
\begin{proof}
  This is the declared model definition of length Readout.
\end{proof}
\begin{definition}[signed Length Readout]
  \label{def:physics:phyx-mini-0610:phyxminiproblems-problemphyxmini0610-signedlengthreadout}
  \lean{PhyXMiniProblems.ProblemPhyXMini0610.signedLengthReadout}
  \uses{def:physics:phyx-mini-0610:phyxminiproblems-problemphyxmini0610-signedlengthquantity}
  Read a signed position or ct coordinate in a selected length unit
\end{definition}
\begin{proof}
  This is the declared model definition of signed Length Readout.
\end{proof}
\begin{definition}[time Readout]
  \label{def:physics:phyx-mini-0610:phyxminiproblems-problemphyxmini0610-timereadout}
  \lean{PhyXMiniProblems.ProblemPhyXMini0610.timeReadout}
  \uses{def:physics:phyx-mini-0610:phyxminiproblems-problemphyxmini0610-timequantity}
  Read a duration in a selected time unit
\end{definition}
\begin{proof}
  This is the declared model definition of time Readout.
\end{proof}
\begin{definition}[speed Readout]
  \label{def:physics:phyx-mini-0610:phyxminiproblems-problemphyxmini0610-speedreadout}
  \lean{PhyXMiniProblems.ProblemPhyXMini0610.speedReadout}
  \uses{def:physics:phyx-mini-0610:phyxminiproblems-problemphyxmini0610-signedspeedquantity}
  Read a signed speed in selected length and time units
\end{definition}
\begin{proof}
  This is the declared model definition of speed Readout.
\end{proof}
\begin{definition}[area Readout]
  \label{def:physics:phyx-mini-0610:phyxminiproblems-problemphyxmini0610-areareadout}
  \lean{PhyXMiniProblems.ProblemPhyXMini0610.areaReadout}
  Read a physical plot area in the square of a selected length unit
\end{definition}
\begin{proof}
  This is the declared model definition of area Readout.
\end{proof}
\begin{definition}[Inertial Frame Label]
  \label{def:physics:phyx-mini-0610:phyxminiproblems-problemphyxmini0610-inertialframelabel}
  \lean{PhyXMiniProblems.ProblemPhyXMini0610.InertialFrameLabel}
  The station frame S and the rocket-train frame S'
\end{definition}
\begin{proof}
  This is the declared model definition of Inertial Frame Label.
\end{proof}
\begin{definition}[Tube Endpoint]
  \label{def:physics:phyx-mini-0610:phyxminiproblems-problemphyxmini0610-tubeendpoint}
  \lean{PhyXMiniProblems.ProblemPhyXMini0610.TubeEndpoint}
  The two physical endpoints of the fluorescent tube
\end{definition}
\begin{proof}
  This is the declared model definition of Tube Endpoint.
\end{proof}
\begin{definition}[Tube Switch Transition]
  \label{def:physics:phyx-mini-0610:phyxminiproblems-problemphyxmini0610-tubeswitchtransition}
  \lean{PhyXMiniProblems.ProblemPhyXMini0610.TubeSwitchTransition}
  The switch transition represented by an endpoint event
\end{definition}
\begin{proof}
  This is the declared model definition of Tube Switch Transition.
\end{proof}
\begin{definition}[Event Label]
  \label{def:physics:phyx-mini-0610:phyxminiproblems-problemphyxmini0610-eventlabel}
  \lean{PhyXMiniProblems.ProblemPhyXMini0610.EventLabel}
  Events numbered 1 through 4 in the supplied spacetime diagram
\end{definition}
\begin{proof}
  This is the declared model definition of Event Label.
\end{proof}
\begin{definition}[event Endpoint]
  \label{def:physics:phyx-mini-0610:phyxminiproblems-problemphyxmini0610-eventendpoint}
  \lean{PhyXMiniProblems.ProblemPhyXMini0610.eventEndpoint}
  \uses{def:physics:phyx-mini-0610:phyxminiproblems-problemphyxmini0610-tubeendpoint, def:physics:phyx-mini-0610:phyxminiproblems-problemphyxmini0610-eventlabel}
  Tube endpoint represented by each numbered event
\end{definition}
\begin{proof}
  This is the declared model definition of event Endpoint.
\end{proof}
\begin{definition}[event Switch Transition]
  \label{def:physics:phyx-mini-0610:phyxminiproblems-problemphyxmini0610-eventswitchtransition}
  \lean{PhyXMiniProblems.ProblemPhyXMini0610.eventSwitchTransition}
  \uses{def:physics:phyx-mini-0610:phyxminiproblems-problemphyxmini0610-tubeswitchtransition, def:physics:phyx-mini-0610:phyxminiproblems-problemphyxmini0610-eventlabel}
  Turning-on or turning-off transition represented by each event
\end{definition}
\begin{proof}
  This is the declared model definition of event Switch Transition.
\end{proof}
\begin{definition}[Spacetime Diagram Point]
  \label{def:physics:phyx-mini-0610:phyxminiproblems-problemphyxmini0610-spacetimediagrampoint}
  \lean{PhyXMiniProblems.ProblemPhyXMini0610.SpacetimeDiagramPoint}
  \uses{def:physics:phyx-mini-0610:phyxminiproblems-problemphyxmini0610-signedlengthquantity}
  A point in an x--ct plot; both coordinates have length dimension
\end{definition}
\begin{proof}
  This is the declared model definition of Spacetime Diagram Point.
\end{proof}
\begin{definition}[Diagram Axis Quantity]
  \label{def:physics:phyx-mini-0610:phyxminiproblems-problemphyxmini0610-diagramaxisquantity}
  \lean{PhyXMiniProblems.ProblemPhyXMini0610.DiagramAxisQuantity}
  The physical quantities represented by the two drawn axes
\end{definition}
\begin{proof}
  This is the declared model definition of Diagram Axis Quantity.
\end{proof}
\begin{definition}[Rectangle Corner]
  \label{def:physics:phyx-mini-0610:phyxminiproblems-problemphyxmini0610-rectanglecorner}
  \lean{PhyXMiniProblems.ProblemPhyXMini0610.RectangleCorner}
  The four corners of the yellow rectangle in image 610
\end{definition}
\begin{proof}
  This is the declared model definition of Rectangle Corner.
\end{proof}
\begin{definition}[Rectangle Edge]
  \label{def:physics:phyx-mini-0610:phyxminiproblems-problemphyxmini0610-rectangleedge}
  \lean{PhyXMiniProblems.ProblemPhyXMini0610.RectangleEdge}
  Boundary segments of the yellow rectangle
\end{definition}
\begin{proof}
  This is the declared model definition of Rectangle Edge.
\end{proof}
\begin{definition}[Boundary Line Style]
  \label{def:physics:phyx-mini-0610:phyxminiproblems-problemphyxmini0610-boundarylinestyle}
  \lean{PhyXMiniProblems.ProblemPhyXMini0610.BoundaryLineStyle}
  Solid and dashed line styles visible on the rectangle boundary
\end{definition}
\begin{proof}
  This is the declared model definition of Boundary Line Style.
\end{proof}
\begin{definition}[Tube Or Region Color]
  \label{def:physics:phyx-mini-0610:phyxminiproblems-problemphyxmini0610-tubeorregioncolor}
  \lean{PhyXMiniProblems.ProblemPhyXMini0610.TubeOrRegionColor}
  Colors relevant to the tube and shaded diagram region
\end{definition}
\begin{proof}
  This is the declared model definition of Tube Or Region Color.
\end{proof}
\begin{definition}[Diagram Text Label]
  \label{def:physics:phyx-mini-0610:phyxminiproblems-problemphyxmini0610-diagramtextlabel}
  \lean{PhyXMiniProblems.ProblemPhyXMini0610.DiagramTextLabel}
  Literal labels visible in the primary raster
\end{definition}
\begin{proof}
  This is the declared model definition of Diagram Text Label.
\end{proof}
\begin{definition}[Axis Direction]
  \label{def:physics:phyx-mini-0610:phyxminiproblems-problemphyxmini0610-axisdirection}
  \lean{PhyXMiniProblems.ProblemPhyXMini0610.AxisDirection}
  Direction of the train's motion along the common spatial axis
\end{definition}
\begin{proof}
  This is the declared model definition of Axis Direction.
\end{proof}
\begin{definition}[Spatial Axis Relationship]
  \label{def:physics:phyx-mini-0610:phyxminiproblems-problemphyxmini0610-spatialaxisrelationship}
  \lean{PhyXMiniProblems.ProblemPhyXMini0610.SpatialAxisRelationship}
  Relationship between the spatial axes of the two frames
\end{definition}
\begin{proof}
  This is the declared model definition of Spatial Axis Relationship.
\end{proof}
\begin{definition}[Fluorescent Tube Spacetime Figure]
  \label{def:physics:phyx-mini-0610:phyxminiproblems-problemphyxmini0610-fluorescenttubespacetimefigure}
  \lean{PhyXMiniProblems.ProblemPhyXMini0610.FluorescentTubeSpacetimeFigure}
  \uses{def:physics:phyx-mini-0610:phyxminiproblems-problemphyxmini0610-eventlabel, def:physics:phyx-mini-0610:phyxminiproblems-problemphyxmini0610-diagramaxisquantity, def:physics:phyx-mini-0610:phyxminiproblems-problemphyxmini0610-rectanglecorner, def:physics:phyx-mini-0610:phyxminiproblems-problemphyxmini0610-rectangleedge, def:physics:phyx-mini-0610:phyxminiproblems-problemphyxmini0610-boundarylinestyle, def:physics:phyx-mini-0610:phyxminiproblems-problemphyxmini0610-tubeorregioncolor, def:physics:phyx-mini-0610:phyxminiproblems-problemphyxmini0610-diagramtextlabel}
  Presentation information transcribed from image 610. The corner assignment, axis meanings, labels, yellow fill, and solid/dashed boundary styles are raw figure evidence; this structure contains no area value or formula
\end{definition}
\begin{proof}
  This is the declared model definition of Fluorescent Tube Spacetime Figure.
\end{proof}
\begin{definition}[Fluorescent Tube Spacetime Setup]
  \label{def:physics:phyx-mini-0610:phyxminiproblems-problemphyxmini0610-fluorescenttubespacetimesetup}
  \lean{PhyXMiniProblems.ProblemPhyXMini0610.FluorescentTubeSpacetimeSetup}
  \uses{def:physics:phyx-mini-0610:phyxminiproblems-problemphyxmini0610-lengthquantity, def:physics:phyx-mini-0610:phyxminiproblems-problemphyxmini0610-timequantity, def:physics:phyx-mini-0610:phyxminiproblems-problemphyxmini0610-signedspeedquantity, def:physics:phyx-mini-0610:phyxminiproblems-problemphyxmini0610-inertialframelabel, def:physics:phyx-mini-0610:phyxminiproblems-problemphyxmini0610-eventlabel, def:physics:phyx-mini-0610:phyxminiproblems-problemphyxmini0610-spacetimediagrampoint, def:physics:phyx-mini-0610:phyxminiproblems-problemphyxmini0610-tubeorregioncolor, def:physics:phyx-mini-0610:phyxminiproblems-problemphyxmini0610-axisdirection, def:physics:phyx-mini-0610:phyxminiproblems-problemphyxmini0610-spatialaxisrelationship, def:physics:phyx-mini-0610:phyxminiproblems-problemphyxmini0610-fluorescenttubespacetimefigure}
  The tube, train, events, illuminated worldsheet, and a general area assignment for regions in the S plot. The area map is not defined from c, L, or T; its required geometric behavior is stated separately below
\end{definition}
\begin{proof}
  This is the declared model definition of Fluorescent Tube Spacetime Setup.
\end{proof}
\begin{definition}[event XReadout]
  \label{def:physics:phyx-mini-0610:phyxminiproblems-problemphyxmini0610-eventxreadout}
  \lean{PhyXMiniProblems.ProblemPhyXMini0610.eventXReadout}
  \uses{def:physics:phyx-mini-0610:phyxminiproblems-problemphyxmini0610-signedlengthreadout, def:physics:phyx-mini-0610:phyxminiproblems-problemphyxmini0610-inertialframelabel, def:physics:phyx-mini-0610:phyxminiproblems-problemphyxmini0610-eventlabel, def:physics:phyx-mini-0610:phyxminiproblems-problemphyxmini0610-fluorescenttubespacetimesetup}
  Spatial-coordinate readout of an event in a chosen frame
\end{definition}
\begin{proof}
  This is the declared model definition of event XReadout.
\end{proof}
\begin{definition}[event Ct Readout]
  \label{def:physics:phyx-mini-0610:phyxminiproblems-problemphyxmini0610-eventctreadout}
  \lean{PhyXMiniProblems.ProblemPhyXMini0610.eventCtReadout}
  \uses{def:physics:phyx-mini-0610:phyxminiproblems-problemphyxmini0610-signedlengthreadout, def:physics:phyx-mini-0610:phyxminiproblems-problemphyxmini0610-inertialframelabel, def:physics:phyx-mini-0610:phyxminiproblems-problemphyxmini0610-eventlabel, def:physics:phyx-mini-0610:phyxminiproblems-problemphyxmini0610-fluorescenttubespacetimesetup}
  ct-coordinate readout of an event in a chosen frame
\end{definition}
\begin{proof}
  This is the declared model definition of event Ct Readout.
\end{proof}
\begin{definition}[point XReadout]
  \label{def:physics:phyx-mini-0610:phyxminiproblems-problemphyxmini0610-pointxreadout}
  \lean{PhyXMiniProblems.ProblemPhyXMini0610.pointXReadout}
  \uses{def:physics:phyx-mini-0610:phyxminiproblems-problemphyxmini0610-signedlengthreadout, def:physics:phyx-mini-0610:phyxminiproblems-problemphyxmini0610-spacetimediagrampoint}
  Spatial-coordinate readout of an arbitrary diagram point
\end{definition}
\begin{proof}
  This is the declared model definition of point XReadout.
\end{proof}
\begin{definition}[point Ct Readout]
  \label{def:physics:phyx-mini-0610:phyxminiproblems-problemphyxmini0610-pointctreadout}
  \lean{PhyXMiniProblems.ProblemPhyXMini0610.pointCtReadout}
  \uses{def:physics:phyx-mini-0610:phyxminiproblems-problemphyxmini0610-signedlengthreadout, def:physics:phyx-mini-0610:phyxminiproblems-problemphyxmini0610-spacetimediagrampoint}
  ct-coordinate readout of an arbitrary diagram point
\end{definition}
\begin{proof}
  This is the declared model definition of point Ct Readout.
\end{proof}
\begin{definition}[Matches Fluorescent Tube Scenario]
  \label{def:physics:phyx-mini-0610:phyxminiproblems-problemphyxmini0610-matchesfluorescenttubescenario}
  \lean{PhyXMiniProblems.ProblemPhyXMini0610.MatchesFluorescentTubeScenario}
  \uses{def:physics:phyx-mini-0610:phyxminiproblems-problemphyxmini0610-fluorescenttubespacetimesetup}
  Qualitative frame, motion, and tube-color assignments from the prose
\end{definition}
\begin{proof}
  This is the declared model definition of Matches Fluorescent Tube Scenario.
\end{proof}
\begin{definition}[Matches Primary Spacetime Figure]
  \label{def:physics:phyx-mini-0610:phyxminiproblems-problemphyxmini0610-matchesprimaryspacetimefigure}
  \lean{PhyXMiniProblems.ProblemPhyXMini0610.MatchesPrimarySpacetimeFigure}
  \uses{def:physics:phyx-mini-0610:phyxminiproblems-problemphyxmini0610-fluorescenttubespacetimesetup}
  Literal figure appearance and the four corner/event assignments. In particular, the primary raster places events 1,2,3,4 at lower-right, lower-left, upper-right, upper-left respectively
\end{definition}
\begin{proof}
  This is the declared model definition of Matches Primary Spacetime Figure.
\end{proof}
\begin{definition}[Matches Event Coordinate Data]
  \label{def:physics:phyx-mini-0610:phyxminiproblems-problemphyxmini0610-matcheseventcoordinatedata}
  \lean{PhyXMiniProblems.ProblemPhyXMini0610.MatchesEventCoordinateData}
  \uses{def:physics:phyx-mini-0610:phyxminiproblems-problemphyxmini0610-lengthreadout, def:physics:phyx-mini-0610:phyxminiproblems-problemphyxmini0610-timereadout, def:physics:phyx-mini-0610:phyxminiproblems-problemphyxmini0610-fluorescenttubespacetimesetup, def:physics:phyx-mini-0610:phyxminiproblems-problemphyxmini0610-eventxreadout, def:physics:phyx-mini-0610:phyxminiproblems-problemphyxmini0610-eventctreadout}
  Calibrated event data from the prose and primary figure. Events 1 and 2 occur at station time zero, while 3 and 4 occur at T. The right-end events have x = 0, and the left-end events have x = -L. Event 1 is also the origin of the primed diagram, as stipulated in the problem
\end{definition}
\begin{proof}
  This is the declared model definition of Matches Event Coordinate Data.
\end{proof}
\begin{definition}[Has Physical Parameters]
  \label{def:physics:phyx-mini-0610:phyxminiproblems-problemphyxmini0610-hasphysicalparameters}
  \lean{PhyXMiniProblems.ProblemPhyXMini0610.HasPhysicalParameters}
  \uses{def:physics:phyx-mini-0610:phyxminiproblems-problemphyxmini0610-lengthreadout, def:physics:phyx-mini-0610:phyxminiproblems-problemphyxmini0610-timereadout, def:physics:phyx-mini-0610:phyxminiproblems-problemphyxmini0610-speedreadout, def:physics:phyx-mini-0610:phyxminiproblems-problemphyxmini0610-fluorescenttubespacetimesetup}
  Positivity and subluminality conditions on the physical parameters
\end{definition}
\begin{proof}
  This is the declared model definition of Has Physical Parameters.
\end{proof}
\begin{definition}[Satisfies Light Scaled Time Coordinate]
  \label{def:physics:phyx-mini-0610:phyxminiproblems-problemphyxmini0610-satisfieslightscaledtimecoordinate}
  \lean{PhyXMiniProblems.ProblemPhyXMini0610.SatisfiesLightScaledTimeCoordinate}
  \uses{def:physics:phyx-mini-0610:phyxminiproblems-problemphyxmini0610-timereadout, def:physics:phyx-mini-0610:phyxminiproblems-problemphyxmini0610-speedreadout, def:physics:phyx-mini-0610:phyxminiproblems-problemphyxmini0610-eventlabel, def:physics:phyx-mini-0610:phyxminiproblems-problemphyxmini0610-fluorescenttubespacetimesetup, def:physics:phyx-mini-0610:phyxminiproblems-problemphyxmini0610-eventctreadout}
  The vertical coordinate convention ct = c t, stated for every event and compatible pair of named units. This is the governing coordinate law, not an assumption about the yellow region's area
\end{definition}
\begin{proof}
  This is the declared model definition of Satisfies Light Scaled Time Coordinate.
\end{proof}
\begin{definition}[Satisfies Illuminated Tube Worldsheet]
  \label{def:physics:phyx-mini-0610:phyxminiproblems-problemphyxmini0610-satisfiesilluminatedtubeworldsheet}
  \lean{PhyXMiniProblems.ProblemPhyXMini0610.SatisfiesIlluminatedTubeWorldsheet}
  \uses{def:physics:phyx-mini-0610:phyxminiproblems-problemphyxmini0610-fluorescenttubespacetimesetup, def:physics:phyx-mini-0610:phyxminiproblems-problemphyxmini0610-eventxreadout, def:physics:phyx-mini-0610:phyxminiproblems-problemphyxmini0610-eventctreadout, def:physics:phyx-mini-0610:phyxminiproblems-problemphyxmini0610-pointxreadout, def:physics:phyx-mini-0610:phyxminiproblems-problemphyxmini0610-pointctreadout}
  The yellow set is exactly the set of lit-tube events and, in the station plot, has the axis-aligned bounds determined by events 2, 1, and 3. These membership facts specify the illuminated worldsheet without assigning it the requested area
\end{definition}
\begin{proof}
  This is the declared model definition of Satisfies Illuminated Tube Worldsheet.
\end{proof}
\begin{definition}[Satisfies Axis Aligned Rectangle Area Law]
  \label{def:physics:phyx-mini-0610:phyxminiproblems-problemphyxmini0610-satisfiesaxisalignedrectanglearealaw}
  \lean{PhyXMiniProblems.ProblemPhyXMini0610.SatisfiesAxisAlignedRectangleAreaLaw}
  \uses{def:physics:phyx-mini-0610:phyxminiproblems-problemphyxmini0610-areareadout, def:physics:phyx-mini-0610:phyxminiproblems-problemphyxmini0610-fluorescenttubespacetimesetup, def:physics:phyx-mini-0610:phyxminiproblems-problemphyxmini0610-eventxreadout, def:physics:phyx-mini-0610:phyxminiproblems-problemphyxmini0610-eventctreadout}
  The ordinary area law for the axis-aligned rectangle bounded by the numbered events. It relates the region's area to its coordinate width and height; it does not contain the requested expression c L T
\end{definition}
\begin{proof}
  This is the declared model definition of Satisfies Axis Aligned Rectangle Area Law.
\end{proof}
\begin{definition}[Answer Choice]
  \label{def:physics:phyx-mini-0610:phyxminiproblems-problemphyxmini0610-answerchoice}
  \lean{PhyXMiniProblems.ProblemPhyXMini0610.AnswerChoice}
  Labels of the four answer choices in the source
\end{definition}
\begin{proof}
  This is the declared model definition of Answer Choice.
\end{proof}
\begin{definition}[printed Answer Expression]
  \label{def:physics:phyx-mini-0610:phyxminiproblems-problemphyxmini0610-printedanswerexpression}
  \lean{PhyXMiniProblems.ProblemPhyXMini0610.printedAnswerExpression}
  \uses{def:physics:phyx-mini-0610:phyxminiproblems-problemphyxmini0610-answerchoice}
  The literal symbolic expression printed beside each choice
\end{definition}
\begin{proof}
  This is the declared model definition of printed Answer Expression.
\end{proof}
\begin{definition}[recorded Dataset Answer]
  \label{def:physics:phyx-mini-0610:phyxminiproblems-problemphyxmini0610-recordeddatasetanswer}
  \lean{PhyXMiniProblems.ProblemPhyXMini0610.recordedDatasetAnswer}
  \uses{def:physics:phyx-mini-0610:phyxminiproblems-problemphyxmini0610-answerchoice}
  Answer label recorded by the source dataset; it is not a premise
\end{definition}
\begin{proof}
  This is the declared model definition of recorded Dataset Answer.
\end{proof}
% --- Archon declaration topology end ---
% --- Archon physics formalization source end ---
